\documentclass[11pt]{article}
\usepackage[margin=1in]{geometry}
\usepackage{amsmath,amssymb,amsthm}
\usepackage{boldtensors}
% boldtensors declares the symbol font `cmbboard` from the bbold family, which
% is not present in every TeX installation (it aborts at the start of math
% mode). Redirect it to the always-available AMS `msb` family; the `"`-syntax
% that would actually render a cmbboard glyph is unused here, so this only
% removes the missing-font dependency.
\DeclareSymbolFont{cmbboard}{U}{msb}{m}{n}
\usepackage{hyperref}
\usepackage{booktabs}
\usepackage{enumitem}

\newcommand{\dd}{\mathrm{d}}
\newcommand{\grad}{\nabla}
\newcommand{\Real}{\mathbb{R}}
\newcommand{\norm}[1]{\lVert #1 \rVert}
\newcommand{\code}[1]{\texttt{#1}}
% Long monospace file paths cannot hyphenate; give TeX room to avoid margin
% overruns and let paths break after slashes where marked with \allowbreak.
\emergencystretch=3em

\title{Design: Energetic Mesh Smoothing on General Geometries\\
via Analytic Surface Constraints}
\author{Alejandro Mota}

\begin{document}
\maketitle

\begin{abstract}
Energetic mesh smoothing (EMS) drives each element toward an ideal reference
shape by solving a quasi-static problem in which nodes relax to minimize a
stored-energy functional. Interior nodes move freely; boundary nodes must
remain on the domain boundary, sliding within it. Today this boundary
restriction is imposed with axis-aligned Dirichlet conditions on the faces of a
cube, which works only because a cube face's outward normal coincides with a
Cartesian basis vector. This note generalizes the boundary restriction to any
boundary that can be described analytically as a level set $g(~x)=0$, using the
same symbolic-expression pipeline already used for initial and boundary
conditions. It states the constraint model (tangential slip with a
return-to-surface correction, and a dimensional hierarchy for faces, edges, and
vertices), two enforcement routes (a soft penalty and an exact local-frame
constraint), the robustness guards, and a test plan that includes exact
reduction to the current cube behavior.
\end{abstract}

\section{Problem statement}

EMS treats the mesh as a hyperelastic body. For each element a target
(``ideal'') reference configuration is constructed from a size criterion
(equal-volume, average-edge-length, \code{max}, or a user size field; see
\code{create\_smooth\_reference} in \code{model.jl}), and a quasi-static solve
lets the nodal positions $~x$ relax so that elements adopt shapes as close as
possible to their targets, minimizing a stored-energy functional $W(~x)$.

Interior nodes are free in $\Real^3$. Boundary nodes are \emph{not}: a node on
the domain boundary $\partial\Omega$ must stay on $\partial\Omega$, or the
smoothing would change the geometry. Physically:
\begin{itemize}[nosep]
  \item a node on a boundary \emph{face} may slide anywhere on that face
        (2 degrees of freedom);
  \item a node on a boundary \emph{edge} (where two faces meet) may slide only
        along the edge (1 degree of freedom);
  \item a node on a boundary \emph{vertex} must stay put (0 degrees of freedom).
\end{itemize}

For a cube, all of this is imposed trivially. Each face is axis-aligned, so
``stay on the face $x=\text{const}$'' is exactly ``hold the $x$ component
fixed, leave $y,z$ free.'' The current EMS examples do precisely that with six
Dirichlet conditions, one Cartesian component per face
(\code{examples/\allowbreak ems/\allowbreak awful-cube/\allowbreak awful-cube.yaml}):
edges and vertices need no special treatment because a node belonging to two or
three faces inherits two or three component constraints automatically.

This is a coincidence of alignment. The moment the boundary is not
axis-aligned---a cylinder's lateral surface, a sphere, a fillet---there is no
Cartesian component to fix, and the current mechanism cannot express the
constraint at all. The goal of this note is to lift EMS to any boundary whose
faces admit an analytic description, so that the method can be explored on
realistic geometries.

\section{Geometric constraint model}
\label{sec:model}

\subsection{A boundary face as a level set}

Assume each smooth boundary face is given as the zero level set of a scalar
field,
\begin{equation}
  g(~x) = 0, \qquad ~x \in \Real^3,
  \label{eq:levelset}
\end{equation}
with $g$ smooth and $\grad g \ne ~0$ on the face (a regular level set). The
outward normal direction is
\begin{equation}
  ~n(~x) = \frac{\grad g(~x)}{\norm{\grad g(~x)}}.
  \label{eq:normal}
\end{equation}
A node constrained to this face may move only tangentially. To first order in
an increment $\delta~x$,
\begin{equation}
  g(~x + \delta~x) \approx g(~x) + \grad g(~x)\cdot\delta~x,
\end{equation}
so requiring the node to remain on the face, $g(~x+\delta~x)=0$, gives the
linearized constraint
\begin{equation}
  \grad g(~x)\cdot\delta~x \;=\; -\,g(~x).
  \label{eq:linconstraint}
\end{equation}
The right-hand side is a \emph{return-to-surface} (drift) correction: it
vanishes when the node lies exactly on the face and otherwise pulls it back
onto the face as the quasi-static iteration proceeds. On a curved face, where a
tangential step necessarily leaves the surface at second order, this term is
what keeps the node on the true geometry rather than on the tangent plane.

Dividing by $\norm{\grad g}$ writes the same constraint in terms of the unit
normal and the signed distance-like residual,
\begin{equation}
  ~n(~x)\cdot\delta~x \;=\; -\,\frac{g(~x)}{\norm{\grad g(~x)}}.
  \label{eq:normalconstraint}
\end{equation}
When the face is axis-aligned, $\grad g = \pm~e_k$, and
\eqref{eq:normalconstraint} reduces to ``fix the $k$-th Cartesian
component''---exactly the present cube behavior. The general case is the
classical \emph{inclined (skew) roller support}: the node is free to slide in
the plane tangent to the surface and pinned in the normal direction.

\subsection{Edges and vertices: the dimensional hierarchy}

A boundary edge is the intersection of two faces, and a vertex the intersection
of three. A node is therefore constrained by \emph{all} the faces it belongs
to. If it satisfies $m$ independent level sets $g_1,\dots,g_m$ (with
$m\in\{1,2,3\}$), it must satisfy \eqref{eq:linconstraint} for each:
\begin{equation}
  \grad g_i(~x)\cdot\delta~x = -\,g_i(~x), \qquad i = 1,\dots,m.
  \label{eq:multiconstraint}
\end{equation}
The constraint gradients $~a_i = \grad g_i(~x)$ span the \emph{constrained
subspace} $\mathcal{C} = \mathrm{span}\{~a_1,\dots,~a_m\}$, of dimension $m$;
the admissible motions are its orthogonal complement, the \emph{tangent
subspace} $\mathcal{T} = \mathcal{C}^\perp$, of dimension $3-m$:

\begin{center}
\begin{tabular}{lccl}
\toprule
feature & level sets $m$ & free DOF $3-m$ & admissible motion \\
\midrule
face   & 1 & 2 & slide in the tangent plane \\
edge   & 2 & 1 & slide along the edge tangent \\
vertex & 3 & 0 & pinned \\
\bottomrule
\end{tabular}
\end{center}

This recovers, with no special cases, the cube behavior in which membership in
one/two/three faces leaves two/one/zero free components. The user declares
which level sets apply to which boundary node set; the feature dimension is
simply the number of level sets listed (or an explicit
\code{surface}/\code{curve}/\code{vertex} tag).

\section{Analytic input and exact gradients}
\label{sec:input}

Norma already compiles analytic expressions of $(t,x,y,z)$ for initial and
boundary conditions and for the EMS size field: the string is parsed with
\code{Meta.parse}, turned into a \code{Symbolics} expression, and compiled to a
fast closure by \code{build\_function} (see
\code{SolidMechanicsDirichletBoundaryCondition} in \code{boundary\_conditions.jl}
and \code{create\_size\_field} in \code{model.jl}). The surface field $g$ is
compiled by the identical path.

The key economy is that the gradient need not be supplied or approximated. From
the symbolic $g$,
\begin{equation}
  \grad g = \code{Symbolics.gradient}(g,\,[x,y,z])
\end{equation}
is obtained \emph{exactly and automatically} and compiled the same way; the
Hessian $\grad^2 g$ is available identically if a full penalty tangent is
wanted. The user writes only the surface, e.g.\ \code{"x\^{}2 + y\^{}2 - R\^{}2"}
for a cylinder of radius $R$ about the $z$-axis; the normal field follows
without finite differencing.

\paragraph{Proposed YAML schema.} Mirroring the initial/boundary-condition
style, each constrained node set lists one or more surfaces; the count sets the
feature dimension.
\begin{verbatim}
model:
  type: mesh smoothing
  smooth reference: max
  boundary surfaces:
    - node set: lateral          # face: 1 level set, 2 free DOF
      surfaces: ["x^2 + y^2 - R^2"]
    - node set: top_rim          # edge: 2 level sets, 1 free DOF
      surfaces: ["x^2 + y^2 - R^2", "z - H"]
    - node set: mount_corner     # vertex: 3 level sets, pinned
      surfaces: ["x", "y", "z - H"]
\end{verbatim}
Constants such as \code{R} and \code{H} are substituted into the string as they
already are for other expressions. A node set with no listed surface is left
free (interior-like), and any Cartesian face can still be written as a linear
level set (e.g.\ \code{"x"}), so the cube case is a special case of the schema.

\section{Enforcement}
\label{sec:enforcement}

The EMS examples solve with matrix-free descent (steepest descent or L-BFGS;
see \code{examples/\allowbreak ems/\allowbreak awful-cube/}), which uses the gradient of $W$ and, for
Newton, its Hessian. Both enforcement routes below fit that machinery; they
differ in invasiveness and in how exactly the surface is honored. Recall the
solver reduces the system to the free DOF: for Newton,
$~A = \code{hessian}[\text{free},\text{free}]$,
$~b = \code{gradient}[\text{free}]$, step $= -~A^{-1}~b$; for matrix-free
descent, the step is taken along the (projected) negative gradient over the
free DOF (\code{solver.jl}).

\subsection{Route A --- penalty (prototype)}
\label{sec:penalty}

Augment the smoothing energy with a quadratic penalty on each constrained
node's surface residuals,
\begin{equation}
  P(~x) = \sum_{\text{nodes}} \; \sum_{i=1}^{m}
          \tfrac12\,\kappa_i\, g_i(~x)^2 .
  \label{eq:penalty}
\end{equation}
Its contributions to the residual (gradient) and tangent (Hessian) at a
constrained node are
\begin{align}
  \frac{\partial P}{\partial ~x}
    &= \sum_i \kappa_i\, g_i\, \grad g_i,
    \label{eq:penaltygrad}\\
  \frac{\partial^2 P}{\partial ~x^2}
    &= \sum_i \kappa_i\bigl(\grad g_i\,\grad g_i^{\top}
       + g_i\,\grad^2 g_i\bigr)
     \;\approx\; \sum_i \kappa_i\,\grad g_i\,\grad g_i^{\top},
    \label{eq:penaltyhess}
\end{align}
the last a Gauss--Newton approximation (drop the $g_i\grad^2 g_i$ term, which
vanishes as $g_i\to0$). Each term is a rank-one operator
$\kappa_i\,\grad g_i\grad g_i^{\top}$ that penalizes motion \emph{along the
normal} while leaving tangential motion governed by $W$---a soft version of the
exact roller. For the matrix-free solvers EMS actually uses, only
\eqref{eq:penaltygrad} is needed; it is added to \code{solver.gradient} at the
node's three DOF after \code{evaluate}. No \code{free\_dofs} manipulation and no
change of basis are required.

\emph{Trade-offs.} Trivial to implement and the fastest way to exercise the
level-set and gradient plumbing on a curved mesh. But the surface is satisfied
only approximately---the residual scales like (smoothing force)$/\kappa$---and
large $\kappa$ ill-conditions the system. $\kappa_i$ must be scaled to the
element stiffness and mesh size (e.g.\ $\kappa \sim c\,\mu/h^2$). It is a
derisking and validation tool, not the production mechanism.

\subsection{Route B --- exact local-frame (inclined roller)}
\label{sec:exact}

Enforce \eqref{eq:multiconstraint} exactly by a change of basis at each
constrained node. Orthonormalize the constraint gradients
$\{~a_1,\dots,~a_m\}$ (Gram--Schmidt / thin QR) into $~q_1,\dots,~q_m$ spanning
$\mathcal{C}$, and complete them to an orthonormal frame
\begin{equation}
  ~Q = [\,\underbrace{~q_1 \cdots ~q_m}_{\text{normal}}
        \mid \underbrace{~q_{m+1}\cdots ~q_3}_{\text{tangent}}\,]
        \in \Real^{3\times3},
\end{equation}
whose first $m$ columns span the constrained (normal) directions and last $3-m$
span the tangent. In the rotated coordinates $\widetilde{\delta~x} = ~Q^\top
\delta~x$, the tangential components are free and the normal components are
prescribed by the drift correction. For a single face ($m=1$),
$\widetilde{\delta x}_1 = -g_1/\norm{~a_1}$ from \eqref{eq:normalconstraint};
for $m>1$ the $m$ normal components follow from the small linear system
$~A_{\mathcal C}\,\widetilde{\delta~x}_{1:m} = -~g$ with
$(~A_{\mathcal C})_{ij} = ~a_i\cdot~q_j$.

\paragraph{Two ways to carry this through the solver.}
\begin{itemize}
  \item \emph{Matrix-free descent (the EMS default).} The exact constraint is a
  \emph{projected gradient with return to surface}: (i) project the node's
  gradient onto the tangent subspace, $~g^{\mathrm{tan}} = ( ~I - \sum_i
  ~q_i~q_i^\top)\,~g^{\mathrm{node}}$, so the descent direction has no normal
  component; (ii) after the step, restore the node to the surface by a few
  Newton iterations of $g_i(~x)=0$ along the normal (a closest-point
  projection). This keeps every constrained node on the manifold to tolerance
  at each iteration and needs no assembled operator---the least invasive exact
  route and the natural fit for steepest descent / L-BFGS.
  \item \emph{Newton (\code{HessianMinimizer}).} Apply the nodal rotation to the
  assembled system: left-multiply the constrained node's rows of the residual
  and Hessian by $~Q^\top$ and right-multiply its columns by $~Q$ (a congruence
  transform that preserves symmetry), mark the $m$ normal DOF as fixed in
  \code{free\_dofs} with prescribed increment equal to the drift correction,
  leave $3-m$ tangential DOF free, solve the reduced system, and rotate the
  increment back, $\delta~x = ~Q\,\widetilde{\delta~x}$.
\end{itemize}

Both variants reduce \emph{exactly} to the present behavior when the faces are
axis-aligned: then $~Q$ is a signed permutation, the projection zeroes a
Cartesian component, and the fixed normal DOF is a Cartesian component---the
cube DBCs, recovered as a special case. This is the production mechanism:
variationally consistent, surface-exact up to the return-to-surface tolerance,
and geometry-agnostic.

\section{Robustness and degeneracy guards}
\label{sec:robustness}

\begin{description}[style=unboxed,leftmargin=0pt,itemsep=2pt]
  \item[Vanishing gradient.] $\norm{\grad g}\to0$ marks a critical point of the
  level set (e.g.\ the center of a sphere), where the normal
  \eqref{eq:normal} is undefined. Abort with a diagnostic, consistent with the
  degeneracy policy already adopted for the overlap-energy blending (reject, do
  not regularize).
  \item[Node off the declared surface.] The node set is asserted by the user to
  lie on the surface. If $|g(~x_0)|$ exceeds a tolerance scaled by the local
  element size at setup, warn (or abort): the mesh does not match the declared
  geometry, which is almost always an input error.
  \item[Rank-deficient feature.] At an edge or vertex the gradients
  $\{~a_i\}$ must be linearly independent; two nearly parallel surfaces make the
  frame system $~A_{\mathcal C}$ ill-conditioned (a near-tangential
  intersection). Detect via the smallest singular value and abort as a
  degenerate feature.
  \item[Return-to-surface convergence.] On strongly curved faces the drift
  correction is only first order; cap the per-iteration normal correction and,
  if the closest-point Newton fails to converge, fall back to a bisection along
  the normal so a node never leaves its surface.
\end{description}

\section{Test plan}
\label{sec:tests}

\begin{enumerate}[nosep]
  \item \textbf{Reduction to the cube (regression).} Re-express the six cube
  faces of \code{awful-cube} as linear level sets ($x$, $x-L$, $y$, \dots) and
  verify EMS produces the same smoothed mesh (to solver tolerance) as the
  current axis-aligned Dirichlet conditions. This pins backward compatibility:
  the general path must contain the special one.
  \item \textbf{Cylinder (the motivating case).} Use an existing cylindrical
  mesh (e.g.\ \code{examples/\allowbreak ahead/\allowbreak \dots/torsion/\allowbreak torsion.g}), perturb the
  interior nodes, constrain the lateral boundary to
  $g=x^2+y^2-R^2=0$, and check that (a) the smoothing energy decreases
  monotonically, (b) every boundary node satisfies $|g|<\text{tol}$ after the
  solve, and (c) boundary nodes moved tangentially (the surface did not
  shrink or bulge). Include the top/bottom rims as edges
  ($\{x^2+y^2-R^2,\; z\mp H\}$) to exercise $m=2$.
  \item \textbf{Penalty vs.\ exact.} On the cylinder, confirm the penalty
  surface residual scales as $1/\kappa$ while the exact route holds
  $|g|$ at the return-to-surface tolerance independent of any penalty
  parameter.
  \item \textbf{Sphere / quarter-annulus.} A doubly curved face
  ($x^2+y^2+z^2-R^2$) and a mixed straight/curved boundary to check that the
  normal field and the dimensional hierarchy behave at genuinely
  non-axis-aligned edges and vertices.
\end{enumerate}

\section{Recommendation and phasing}
\label{sec:phasing}

\begin{description}[style=unboxed,leftmargin=0pt,itemsep=2pt]
  \item[Phase 1.] Level-set input (Section~\ref{sec:input}) plus the penalty
  enforcement (Section~\ref{sec:penalty}) and the cylinder test. This
  derisks the input parsing, the \code{Symbolics} gradient path, and the
  membership/normal evaluation with the smallest possible code change, and
  produces visible ``nodes hug the surface'' behavior early.
  \item[Phase 2.] The exact local-frame constraint
  (Section~\ref{sec:exact})---projected-gradient/return-to-surface for the
  matrix-free solvers first, then the Newton nodal rotation---and the
  reduce-to-cube regression. This is the production mechanism.
  \item[Phase 3 (optional).] Edges and vertices ($m>1$) end-to-end, plus
  refinements: a symmetric smoothstep or harmonic weighting if a $C^1$ ramp is
  ever wanted, and caching of the compiled $g,\grad g$ per node set alongside
  the existing size-field closure.
\end{description}

The construction reuses machinery already in the tree---the symbolic-expression
compiler, \code{Symbolics.gradient}, the \code{free\_dofs} reduced solve, and
the \code{compute\_rotation\_matrix} helper in
\code{boundary\_conditions.jl}---and degrades exactly to the current cube
behavior, so the present EMS examples remain a regression test of the general
path.

\end{document}
